\documentclass[book.tex]{subfiles}
\begin{document}

\chapter{Spatial Graph Convolutions}
\label{chap:spatial_graph_convs}
\noindent\rule{11cm}{0.4pt} \\
\textbf{Prerequisites:} \autoref{chap:graphconvs}, \autoref{chap:equivariant_networks}, \autoref{chap:molecular_dynamics},  \autoref{chap:dft} \\
\textbf{Difficulty Level:} ** \\
\noindent\rule{11cm}{0.4pt}
\vspace{5mm}

Spatial graph convolutions use the three dimensional structure of the molecule in the computed representation.

"The arc of drug discovery entails a multiparameter optimization problem spanning vast length scales. Through feature learning—instead of feature engineering—deep neural networks promise to outperform both traditional physics-based and knowledge-based machine learning models for predicting molecular properties pertinent to drug discovery. To this end, we present the PotentialNet family of graph convolutions. Alongside, predicting absorption, distribution, metabolism, elimination, and toxicity (ADMET) properties has therefore been of great interest to the computational chemistry and medicinal chemistry communities in recent decades. Here, we learn the features most relevant to each chemical task at hand by representing each molecule explicitly as a graph. By applying graph convolutions to this explicit molecular representation, we achieve, to our knowledge, unprecedented accuracy in prediction of ADMET properties."

References \cite{gomes2017atomic}, \cite{danel2020spatial}, \cite{feinberg2018potentialnet}, \cite{feinberg2020improvement}

\end{document}